\origpage{162}{ROBERTSON AND SEYMOUR}
{\centering
{\Large Graph Minors\par}
{\large XVII. Taming a Vortex\par}\medskip

Neil Robertson$^*$\par
\textit{Department of Mathematics, Ohio State University, 231 West 18th Avenue, Columbus, Ohio}\par
and\par
P. D. Seymour\par
\textit{Department of Mathematics, Princeton University, Princeton, New Jersey 08544-1000, and Bellcore, 445 South Street, Morristown, New Jersey 07960}\par
Received August 14, 1992\par
}
\begin{quote}
The main result of this series serves to reduce several problems about general graphs to problems about graphs which can ``almost'' be drawn in surfaces of bounded genus. In applications of the theorem we usually need to encode such a nearly embedded graph as a hypergraph which can be drawn completely in the surface. The purpose of this paper is to show how to ``tidy up'' near-embeddings to facilitate the encoding procedure.
\end{quote}

\sectionhead{1. \textsc{Introduction}}

In [2] we gave a theorem about the structure of graphs with no minor isomorphic to a fixed graph. (Graphs in this paper are finite, and may have loops or multiple edges. A graph $H$ is a minor of a graph $G$ if $H$ can be obtained from a subgraph of $G$ by contracting edges.) That theorem said that for every graph $H$, every graph with no minor isomorphic to $H$ can be expressed as a tree-structure of ``pieces,'' where each piece is a graph which can be drawn in a surface in which $H$ cannot be drawn, except for a bounded number of vertices and a bounded number of ``local areas of non-planarity'' (or vortices, as we call them). (The bounds here depend on $H$, but not on the graph being decomposed.) In applications each piece is usually then encoded as a hypergraph which can be drawn completely on the surface, with edges labelled from a appropriate quasi-order, and in%
\origfootnote{This research was partially supported by NSF grant DMS 8504054 and was partially performed under a consulting agreement with Bellcore.}
